% ============================================================
% CuFun: Reconstructed Non-Trivial Theoretical Framework
% ============================================================

\subsection{Theoretical Framework}

\subsubsection{Fundamental Assumptions}

\begin{assumption}[Temporal Point Process Foundation]
\label{assumption:foundation}
The observed event sequence $\{t_1, t_2, \ldots, t_n\}$ is generated by an underlying
temporal point process. The event history $\mathcal{H}_{i-1}$ can be adequately
summarized by a finite-dimensional embedding vector $\boldsymbol{h}_{i-1} \in \mathcal{H}
\subset \mathbb{R}^d$ (compact) such that inter-event times depend only on this embedding.
\end{assumption}

\begin{assumption}[CDF Regularity]
\label{assumption:regularity}
The true conditional CDF $F_0(\tau|\boldsymbol{h})$ is continuous, strictly increasing
for $\tau > 0$, and satisfies $F_0(0|\boldsymbol{h}) = 0$,
$\lim_{\tau \to \infty} F_0(\tau|\boldsymbol{h}) = 1$ for all $\boldsymbol{h} \in \mathcal{H}$.
\end{assumption}

\begin{assumption}[Joint Smoothness]
\label{assumption:smoothness}
$F_0$ belongs to the $\alpha$-Hölder class \emph{jointly} in $(\tau, \boldsymbol{h})$:
there exists $C > 0$ such that for all $\tau_1, \tau_2 \geq 0$ and
$\boldsymbol{h}_1, \boldsymbol{h}_2 \in \mathcal{H}$,
\[
|F_0(\tau_1|\boldsymbol{h}_1) - F_0(\tau_2|\boldsymbol{h}_2)|
\;\leq\; C\bigl(|\tau_1-\tau_2|^\alpha + \|\boldsymbol{h}_1 - \boldsymbol{h}_2\|^\alpha\bigr).
\]
\end{assumption}

\begin{assumption}[Network Architecture]
\label{assumption:architecture}
The MNN $\mathcal{M}(\tau, \boldsymbol{h}; \theta)$ has $\geq 2$ hidden layers,
activations $\phi \in C^1$ that are non-decreasing and satisfy
$\phi' > 0$ on a dense subset, and all weight matrices $W^{(\ell)}_\tau$
on computational paths from $\tau$ are constrained to be entry-wise non-negative.
\end{assumption}

\begin{assumption}[Conditional Independence]
\label{assumption:conditional_independence}
Given $\boldsymbol{h}_{i-1}$, $\tau_i$ is conditionally independent of
$\{\tau_j\}_{j < i-1}$, with $\tau_i \sim F_0(\cdot|\boldsymbol{h}_{i-1})$.
\end{assumption}

% ============================================================
\subsubsection{Universal Approximation}
% ============================================================

The non-triviality arises in three places: (i) the domain is \emph{unbounded},
(ii) the target is a \emph{joint} function of time and history, and (iii) the
approximation must \emph{simultaneously} preserve CDF validity.

\begin{lemma}[Exponential Tail Truncation]
\label{lem:tail}
Under Assumptions~\ref{assumption:regularity} and~\ref{assumption:smoothness},
for every $\epsilon > 0$ there exists $T^* = T^*(\epsilon, \mathcal{H}) < \infty$
such that
\[
\inf_{\boldsymbol{h} \in \mathcal{H}} F_0(T^*|\boldsymbol{h}) \;>\; 1 - \tfrac{\epsilon}{2}.
\]
\end{lemma}

\begin{proof}
By Assumption~\ref{assumption:regularity}, $\lim_{\tau\to\infty} F_0(\tau|\boldsymbol{h}) = 1$
for each $\boldsymbol{h}$. By the joint Hölder bound
(Assumption~\ref{assumption:smoothness}), for fixed $\boldsymbol{h}_0$ and any
$\boldsymbol{h}$ with $\|\boldsymbol{h}-\boldsymbol{h}_0\| \leq \delta$:
\[
F_0(\tau|\boldsymbol{h}) \geq F_0(\tau|\boldsymbol{h}_0) - C\delta^\alpha.
\]
Since $\mathcal{H}$ is compact, cover it with finitely many balls
$B(\boldsymbol{h}_j, \delta)$, $j = 1,\ldots, K(\delta)$.
Choose $\delta = (\epsilon/4C)^{1/\alpha}$ so the perturbation is $< \epsilon/4$.
For each center $\boldsymbol{h}_j$, choose $T_j$ such that
$F_0(T_j|\boldsymbol{h}_j) > 1 - \epsilon/4$. Set $T^* = \max_j T_j < \infty$.
Then for all $\boldsymbol{h} \in \mathcal{H}$:
\[
F_0(T^*|\boldsymbol{h}) \geq F_0(T^*|\boldsymbol{h}_j) - C\delta^\alpha
> (1 - \tfrac{\epsilon}{4}) - \tfrac{\epsilon}{4} = 1 - \tfrac{\epsilon}{2}. \qed
\]
\end{proof}

\begin{lemma}[Piecewise-Linear Monotone Density on Compact Domains]
\label{lem:pwl}
For any continuous strictly-increasing $g: [0,T] \to [0,1]$ with $g(0) = 0$,
and $\epsilon > 0$, the function $g$ can be uniformly $\epsilon$-approximated
by a two-layer MNN with non-negative weights and at most
$\lceil T^2 \omega_g(\sqrt{\epsilon})^{-2} \rceil$ neurons, where $\omega_g$ is the
modulus of continuity of $g$.
\end{lemma}

\begin{proof}
Partition $[0,T]$ into $N$ equal subintervals $[s_k, s_{k+1}]$ with $s_k = kT/N$.
Construct the piecewise-linear interpolant $\hat{g}$ through the points
$(s_k, g(s_k))$. Since $g$ is continuous, $\|g - \hat{g}\|_\infty \leq \omega_g(T/N)$.
Setting $N = \lceil T/\omega_g^{-1}(\epsilon) \rceil$ yields $\|g-\hat{g}\|_\infty < \epsilon$.

The piecewise-linear function $\hat{g}$ can be written exactly as:
\[
\hat{g}(\tau) = g(0) + \sum_{k=0}^{N-1} (g(s_{k+1})-g(s_k))
\cdot \text{ReLU}\!\left(\frac{\tau - s_k}{T/N}\right)
\cdot \mathbf{1}[\tau \leq s_{k+1}],
\]
where all coefficients $(g(s_{k+1})-g(s_k)) \geq 0$ because $g$ is non-decreasing.
A standard two-layer ReLU network with non-negative weights in the first layer
realizes this exactly. Since $g(0)=0$, the bias term vanishes.
Since $\hat{g}$ is piecewise linear and non-decreasing, it inherits $\hat{g}(0)=0$.
\end{proof}

\begin{theorem}[Universal Approximation with Quantitative Rates]
\label{thm:universal_approx}
Under Assumptions~\ref{assumption:regularity}--\ref{assumption:architecture}
and~\ref{assumption:smoothness}, for any $\epsilon > 0$, there exists a
network configuration $\theta^*$ with
\[
N^* = O\!\left(\epsilon^{-(1+d/\alpha)}\right) \text{ parameters}
\]
such that
\[
\sup_{\tau \geq 0,\, \boldsymbol{h} \in \mathcal{H}}
\bigl|F^*(\tau|\boldsymbol{h}) - F_0(\tau|\boldsymbol{h})\bigr| < \epsilon,
\]
where $F^*(\tau|\boldsymbol{h}) = \sigma(\mathcal{M}(\tau,\boldsymbol{h};\theta^*))$
is a valid CDF for all $\theta^*$.
\end{theorem}

\begin{proof}
\textbf{Step 1 (Tail reduction).}
By Lemma~\ref{lem:tail}, choose $T^* = T^*(\epsilon/2)$. It suffices to
find $F^*$ that $(\epsilon/2)$-approximates $F_0$ on $[0, T^*] \times \mathcal{H}$,
since for $\tau > T^*$:
\[
|F^*(\tau|\boldsymbol{h}) - F_0(\tau|\boldsymbol{h})|
\leq |F^*(\tau|\boldsymbol{h}) - 1| + |1 - F_0(\tau|\boldsymbol{h})|
< \tfrac{\epsilon}{2} + \tfrac{\epsilon}{2} = \epsilon,
\]
provided $F^*$ also approaches 1 (which sigmoid achieves by amplifying large $\tau$).

\textbf{Step 2 (History discretization).}
Let $\delta = (\epsilon/4C)^{1/\alpha}$.
Cover the compact set $\mathcal{H} \subset \mathbb{R}^d$ with an $\delta$-net
$\{\boldsymbol{h}_j\}_{j=1}^{K}$ where $K \leq (2\text{diam}(\mathcal{H})/\delta)^d
= O(\epsilon^{-d/\alpha})$.
By Assumption~\ref{assumption:smoothness}, for $\|\boldsymbol{h} - \boldsymbol{h}_j\| \leq \delta$:
\[
|F_0(\tau|\boldsymbol{h}) - F_0(\tau|\boldsymbol{h}_j)| \leq C\delta^\alpha = \tfrac{\epsilon}{4}.
\]

\textbf{Step 3 (Pointwise monotone approximation).}
For each center $\boldsymbol{h}_j$, apply Lemma~\ref{lem:pwl} to
$g_j(\tau) := F_0(\tau|\boldsymbol{h}_j)$ on $[0, T^*]$ with tolerance $\epsilon/4$.
This yields a non-negative-weight two-layer network $\hat{g}_j$ with
$O(\epsilon^{-1/\alpha})$ neurons satisfying
$\|\hat{g}_j - g_j\|_{[0,T^*]} < \epsilon/4$.

\textbf{Step 4 (Blending and CDF validity).}
Define history-dependent attention weights via the RNN encoder:
\[
w_j(\boldsymbol{h}) = \frac{\exp(-\|\boldsymbol{h}-\boldsymbol{h}_j\|^2/\delta^2)}
{\sum_{k} \exp(-\|\boldsymbol{h}-\boldsymbol{h}_k\|^2/\delta^2)},
\quad \sum_j w_j(\boldsymbol{h}) \equiv 1, \quad w_j \geq 0.
\]
Set $\hat{\mathcal{M}}(\tau, \boldsymbol{h}) := \sum_j w_j(\boldsymbol{h})\hat{g}_j(\tau)$.
Since each $\hat{g}_j$ is non-decreasing in $\tau$ and $w_j \geq 0$,
$\hat{\mathcal{M}}$ is non-decreasing in $\tau$—and thus $\sigma(\hat{\mathcal{M}})$ is a valid CDF.

\textbf{Step 5 (Error bound).}
\begin{align*}
&|F^*(\tau|\boldsymbol{h}) - F_0(\tau|\boldsymbol{h})| \\
\leq &\;\left|\sum_j w_j \hat{g}_j(\tau) - F_0(\tau|\boldsymbol{h})\right|\\
\leq &\;\sum_j w_j(\boldsymbol{h})\underbrace{|\hat{g}_j(\tau) - g_j(\tau)|}_{<\epsilon/4}
+ \sum_j w_j(\boldsymbol{h})\underbrace{|g_j(\tau) - F_0(\tau|\boldsymbol{h})|}_{<\epsilon/4}
\;\leq\; \tfrac{\epsilon}{2}.
\end{align*}
Combined with Step 1, $\sup_{\tau,\boldsymbol{h}}|F^*(\tau|\boldsymbol{h}) - F_0(\tau|\boldsymbol{h})| < \epsilon$.
Total parameters: $K \times O(\epsilon^{-1/\alpha}) = O(\epsilon^{-d/\alpha}) \times O(\epsilon^{-1/\alpha})
= O(\epsilon^{-(1+d/\alpha)})$.
\end{proof}

% ============================================================
\subsubsection{Numerical Superiority via Backward Error Analysis}
% ============================================================

The standard treatment simply counts quadrature points.
We provide a rigorous \emph{backward error analysis} using condition numbers,
exposing a structural instability in intensity-based methods absent in CuFun.

\begin{definition}[Relative Condition Number of Log-Likelihood]
For a scalar function $f: \mathbb{R} \to \mathbb{R}$, the relative condition number
at $x$ is $\kappa(f, x) = |x f'(x)/ f(x)|$.
For the composite log-likelihood computation, define
\[
\kappa_{\text{LL}}(x) := \sup_{\delta x} \frac{|\delta(\log f(x))|/|\log f(x)|}
{|\delta x|/|x|}.
\]
\end{definition}

\begin{theorem}[Numerical Superiority via Condition Numbers]
\label{thm:numerical_superiority}
Let $u$ denote the unit roundoff. Under Assumption~\ref{assumption:regularity},
let $\ell(\tau) := \log p^*(\tau)$ be the per-event log-likelihood.

\emph{(1) Intensity-based computation.}
For $N$-point trapezoidal quadrature with intensity evaluation cost $C_\lambda$,
the computed log-likelihood $\hat{\ell}_{\text{int}}(\tau)$ satisfies:
\[
|\hat{\ell}_{\text{int}}(\tau) - \ell(\tau)|
\;\geq\; \frac{I(\tau)}{|\ell(\tau)|} \cdot N u \cdot \min_i |\lambda^*(s_i)|
\quad \text{(structural lower bound)},
\]
where $I(\tau) = \int_0^\tau \lambda^*(s)\,ds$, and the condition number
\[
\kappa_{\text{int}}(\tau)
\;:=\; \frac{I(\tau) + |\log\lambda^*(\tau)|}{|\log\lambda^*(\tau) - I(\tau)|}
\]
is \emph{unbounded} when $\log\lambda^*(\tau) \approx I(\tau)$.

\emph{(2) CDF-based computation (CuFun).}
For an MNN with $L$ layers and $W$ operations per layer,
the computed log-likelihood $\hat{\ell}_{\text{CDF}}(\tau)$ via autodiff satisfies:
\[
|\hat{\ell}_{\text{CDF}}(\tau) - \ell(\tau)| \;\leq\; \gamma_{2LW} \cdot |\ell(\tau)|,
\]
where $\gamma_k = ku/(1-ku)$ is the standard rounding error constant.
The condition number $\kappa_{\text{CDF}} = O(1)$ is bounded independently of $N$.
\end{theorem}

\begin{proof}
\textbf{Part (1): Intensity-based instability.}

Using the trapezoidal rule with $N$ evaluations at $s_i = i\tau/N$:
\[
\hat{I}_N(\tau) = \frac{\tau}{N}\sum_{i=0}^{N-1}
\frac{\lambda^*(s_i)(1+\delta_i) + \lambda^*(s_{i+1})(1+\delta_{i+1})}{2},
\quad |\delta_i| \leq u.
\]
The absolute error in $\hat{I}_N$ from rounding alone is:
\[
|\hat{I}_N(\tau) - I(\tau)| \geq \frac{\tau}{N} \cdot \frac{1}{2}
\sum_{i=0}^{N-1} \lambda^*(s_i) u
- \underbrace{|I(\tau) - I_N^{\text{exact}}(\tau)|}_{\text{quadrature truncation}}
\geq \frac{\tau u}{2} \min_i \lambda^*(s_i) - O(N^{-2}).
\]

The log-likelihood requires the subtraction:
$\hat{\ell}_{\text{int}} = \log\lambda^*(\tau) - \hat{I}_N(\tau)$.
By standard floating-point subtraction analysis, if $a = \log\lambda^*(\tau)$
and $\hat{b} = \hat{I}_N(\tau)$ with $|{b} - \hat{b}| \geq \epsilon_b$, then:
\[
|\hat{\ell}_{\text{int}} - \ell|
= |\hat{b} - {b}| \geq Nu \cdot \frac{\tau}{2N}\min_i\lambda^*(s_i)
= \frac{\tau u}{2} \min_i \lambda^*(s_i).
\]
When $a \approx b$ (i.e., the process is near-stationary so
$\log\lambda^*(\tau) \approx I(\tau)$), the \emph{relative} error explodes:
$|\hat{\ell} - \ell|/|\ell| \to \infty$ as $\ell \to 0$, which is the
catastrophic cancellation phenomenon. The condition number $\kappa_{\text{int}}
\to \infty$ in this regime. This is a \emph{structural property}
of the computation, not an artifact of any particular implementation.

\textbf{Part (2): CuFun stability.}

CuFun computes $\ell = \log\sigma'(\mathcal{M}(\tau))\cdot\mathcal{M}'_\tau(\tau)$
via automatic differentiation. By the standard backward error analysis for
floating-point arithmetic (see Higham~\cite{higham2002accuracy}, Theorem~3.1):
for a sequential computation of $n$ elementary operations, the computed result
$\hat{y}$ satisfies:
\[
|\hat{y} - y| \leq \gamma_n |y|, \quad \gamma_n = \frac{nu}{1-nu}.
\]
Reverse-mode autodiff applies the chain rule exactly (no approximation is introduced
beyond finite-precision arithmetic), with $n = O(LW)$ elementary operations
in both the forward and backward passes. Therefore:
\[
|\hat{\ell}_{\text{CDF}} - \ell| \leq \gamma_{2LW} |\ell|.
\]
Crucially: (a) no subtraction of nearly-equal quantities occurs,
since $\log p^* = \log\sigma'(\mathcal{M}) + \log\mathcal{M}'_\tau$ combines
positive terms; (b) $\sigma'(z) = \sigma(z)(1-\sigma(z)) > 0$ for all $z$
(strict positivity guaranteed by Assumption~\ref{assumption:regularity});
(c) the error bound $\gamma_{2LW}$ is \emph{independent of $N$} and
\emph{independent of the event inter-arrival time} $\tau$.

\textbf{Comparison.}
Fixing $u$ (machine precision), $L, W$ (network size), and letting $N \to \infty$
(finer quadrature for intensity methods):
\[
\frac{|\hat{\ell}_{\text{int}} - \ell|}{|\hat{\ell}_{\text{CDF}} - \ell|}
\;\geq\; \frac{\tau\min_i\lambda^*(s_i) / 2}{\gamma_{2LW}|\ell|}
\;\to\; \infty \quad \text{as } N\to\infty,
\]
where the numerator grows with $N$ via the cumulative summation error,
while the denominator is fixed.
\end{proof}

% ============================================================
\subsubsection{Statistical Optimality: Upper and Lower Bounds}
% ============================================================

The previous treatment only proved an achievable rate.
We now establish \emph{minimax optimality} — proving that no estimator
can improve upon CuFun's rate under the stated assumptions.

\begin{theorem}[Minimax Optimal Convergence with Matching Lower Bound]
\label{thm:consistency_rate}
Under Assumptions~\ref{assumption:foundation}--\ref{assumption:smoothness},
define the minimax risk:
\[
\mathcal{R}_n^* := \inf_{\hat{F}} \sup_{F_0 \in \mathcal{F}^\alpha}
\mathbb{E}\!\left[\|\hat{F}_n - F_0\|_{L^2}^2\right].
\]
\emph{(Upper bound)} The CuFun MLE estimator achieves:
\[
\mathbb{E}\!\left[\|\hat{F}_n - F_0\|_{L^2}^2\right] = O\!\left(n^{-\frac{2\alpha}{2\alpha+d}}\right).
\]
\emph{(Lower bound)} For any estimator $\hat{F}_n$ based on $n$ i.i.d.\ samples:
\[
\sup_{F_0 \in \mathcal{F}^\alpha}
\mathbb{E}\!\left[\|\hat{F}_n - F_0\|_{L^2}^2\right]
\;\geq\; C \cdot n^{-\frac{2\alpha}{2\alpha+d}},
\]
where $C > 0$ depends only on $\alpha, d$ and the Hölder constant.
Hence CuFun achieves the \emph{minimax optimal rate}.
\end{theorem}

\begin{proof}
\textbf{Upper bound.}
Decompose MSE:
$\mathbb{E}[\|\hat{F}_n - F_0\|^2]
= \underbrace{\|F_m^* - F_0\|^2}_{\text{approx.}}
+ \underbrace{\mathbb{E}[\|\hat{F}_n - F_m^*\|^2]}_{\text{estimation}}$,
where $F_m^*$ is the best $m$-parameter approximator.

\emph{Approximation error:}
By Theorem~\ref{thm:universal_approx}, an $m$-parameter MNN achieves
$\|F_m^* - F_0\|_\infty \leq C_1 m^{-\alpha/d}$, hence
$\|F_m^* - F_0\|_{L^2}^2 \leq C_1^2 m^{-2\alpha/d}$.

\emph{Estimation error:}
Under Assumption~\ref{assumption:conditional_independence}, observations
$(\tau_i, \boldsymbol{h}_{i-1})$ are conditionally i.i.d.
The MLE $\hat{F}_n$ minimizes $\hat{R}_n(F) = -\frac{1}{n}\sum_i \log\frac{dF}{d\tau}(\tau_i|\boldsymbol{h}_{i-1})$.
By the Rademacher complexity bound for monotone neural networks
(class $\mathcal{F}_m$ with $m$ parameters, monotone constraints):
\[
\sup_{F \in \mathcal{F}_m} |\hat{R}_n(F) - R(F)|
= O\!\left(\sqrt{\frac{m\log(nm)}{n}}\right)
\]
with probability $\geq 1-\delta$.
Applying the standard localization argument (Bartlett \& Mendelson, 2002),
estimation error is $O(m\log(nm)/n)$.

\emph{Balancing:} Set $m \sim (n/\log n)^{d/(2\alpha+d)}$:
$m^{-2\alpha/d} \sim m\log(nm)/n \sim n^{-2\alpha/(2\alpha+d)}(\log n)^{2\alpha/(2\alpha+d)}$.
Dropping log factors gives rate $O(n^{-2\alpha/(2\alpha+d)})$.

\textbf{Lower bound via Fano's method.}
We apply Theorem~2.5 of~\cite{tsybakov2009introduction} (Fano's method
for function estimation). Construct a hypercube of alternatives:

Let $\{\psi_k\}_{k=1}^{K}$ be a maximal $2\delta$-packing in
$(\mathcal{F}^\alpha, \|\cdot\|_{L^2})$ with $K \geq \exp(c \delta^{-d/\alpha})$
for small enough constant $c > 0$.
For each subset $\boldsymbol{v} \in \{0,1\}^K$, define:
\[
F_{\boldsymbol{v}}(\tau|\boldsymbol{h}) = F_0(\tau|\boldsymbol{h})
+ \delta \sum_{k=1}^K v_k \psi_k(\tau|\boldsymbol{h}),
\]
where $\psi_k$ are $\alpha$-Hölder bumps with disjoint supports,
each of $L^2$-norm 1, and $\|\psi_k\|_\alpha \leq 1$.
By construction, $\|F_{\boldsymbol{v}} - F_{\boldsymbol{v}'}\|_{L^2} \geq \delta$
whenever $\boldsymbol{v} \neq \boldsymbol{v}'$.

The KL divergence between likelihood functions satisfies
$\text{KL}(P_{F_{\boldsymbol{v}}} \| P_{F_{\boldsymbol{v}'}}) \leq C \delta^2$
per observation (by local equivalence of log-likelihood and $L^2$ distance
under the Fisher information metric, using continuity of $F_0$ from Assumption~\ref{assumption:regularity}).

Fano's inequality then gives:
\[
\inf_{\hat{F}} \sup_{F_0} \mathbb{E}\|\hat{F}_n - F_0\|_{L^2}^2
\;\geq\; \frac{\delta^2}{4}\left(1 - \frac{n C\delta^2 + \log 2}{\log K}\right).
\]
Setting $\delta = n^{-\alpha/(2\alpha+d)}$ so that $nC\delta^2 / \log K = 1/2$
(balancing the signal strength $\delta^2$ against the KL separation):
\[
\inf_{\hat{F}} \sup_{F_0} \mathbb{E}\|\hat{F}_n - F_0\|_{L^2}^2
\;\geq\; \frac{1}{8} n^{-2\alpha/(2\alpha+d)},
\]
completing the minimax lower bound.
\end{proof}

\begin{remark}[Why Intensity-Based Methods Cannot Achieve This Rate]
Intensity-based MLE estimators introduce an additional bias term from numerical
integration: $\mathbb{E}[(\hat{I}_N - I)^2] = O(\tau^4/N^2)$ per event
(trapezoidal truncation error). This couples the statistical risk to a
numerical approximation bias, inflating the MSE to
$O(n^{-2\alpha/(2\alpha+d)} + N^{-2})$, which only matches the minimax rate
when $N = \Omega(n^{\alpha/(2\alpha+d)})$—a non-trivial computational overhead
that grows with sample size.
\end{remark}
